\appendix
\section{附录}
\label{sec:appendix}
\subsection{主要运行命令}
在项目根目录安装运行依赖，按功能目录放置模型资产，再启动桌面端。GPU 环境使用 CUDA 参数，CPU 可用于图片和小样本检查。
\begin{lstlisting}[language={},basicstyle=\ttfamily\small,breaklines=true,frame=none]
python -m pip install -r requirements/training.txt
python -m pip install -r requirements/anonymization.txt
python -m pip install -r requirements/ui.txt
python -m expression_deid.desktop_ui --device cuda
\end{lstlisting}

\subsection{模型与工程检查}
运行资源分别放入 \path{models/detection/}、\path{models/anonymization/} 和 \path{models/expression/}，模型清单见表\ref{tab:models}。质量表情权重同时配有 JSON 契约，记录输入预处理、类别顺序、校准参数和校验值。更换权重时应同步核对该文件。
\begin{table}[htbp]
\centering
\caption{桌面运行资源}
\label{tab:models}
\begin{tabularx}{\linewidth}{@{}lX@{}}
\toprule
资源文件 & 用途\\
\midrule
\path{anonymizer_adversarial_v1.pt} & 统一匿名化的局部残差组件\\
\path{ganonymization_50.ckpt} & 关键点条件生成组件\\
\path{face_landmarker.task} & 面部关键点提取\\
\path{face_detection_yunet_2023mar.onnx} & 人脸框与五点检测\\
\path{expression_realtime_stable.pt} & 高速视频表情配置\\
\path{expression_machine_selector_texture.pt} & 图片实验表情配置\\
\path{expression_video_quality.ts.pt} & 质量表情配置\\
\bottomrule
\end{tabularx}
\end{table}

工程检查包括语法与导入检查、匿名输入边界、图像尺寸、页面切换、视频停止、资源释放和异常输入。分类器输入应来自匿名结果，无脸或匿名处理失败时不以源脸补齐。下列入口用于日常检查，逐文件职责与训练产物位置记录于 \path{docs/architecture/project-layout.md}；具体模型效果按数据集评价。
\begin{lstlisting}[language={},basicstyle=\ttfamily\small,breaklines=true,frame=none]
python -m compileall -q expression_deid scripts tests
python -m unittest discover -s tests -p "test_*.py" -q
python scripts/checks/check_desktop_presentation.py --device cuda
\end{lstlisting}

\clearpage
\subsection{第三方资源}
人脸检测使用 YuNet，关键点提取使用 MediaPipe Face Landmarker，条件生成使用 GANonymization 的公开组件。表情质量配置使用 EmotiEffLib 的固定版本表征，其来源和许可记录于 \path{third_party/licenses/expression-quality-NOTICE.md}。训练数据分别遵守 RAF-DB、KMU-FED、PPB-Emo 和 FER+ 的使用要求；工程中保留相应来源，数据与模型按各自许可共享。

报告排版采用 LostSirius 的同济课程报告 LaTeX 模板，运行与界面依赖 PyTorch、OpenCV、PyAV、PySide6。第三方网络和公开权重用于其实际承担的组件，项目工作包括输入链路设计、匿名域适配、局部合成、时序处理、桌面集成与实验分析。

\subsection{项目分工}
\begin{table}[htbp]
\centering
\caption{小组任务分工}
\label{tab:division}
\begin{tabularx}{\linewidth}{@{}lX@{}}
\toprule
成员 & 工作内容\\
\midrule
高霁 & 总体方案与核心链路设计，匿名化与表情模型整合、训练优化及联合验证；桌面端开发与 GPU 推理加速，实验统筹、工程交付、运行文档及报告统稿\\
刘浩 & 数据集资料与许可核对、样本清洗、标签整理、灰度化及人脸对齐与清单检查\\
杨宽 & 表情训练方案设计、参数与日志记录、模型校准及分类别准确率和混淆矩阵分析\\
陈紫坤 & 匿名化文献调研、备选方法比较、局部编辑方案分析及匿名效果与表情保持对照\\
王梓诚 & 视频与摄像头接入、连续推理联调、时序稳定性检查及分析帧率与资源释放记录\\
明子皓 & 实验图表与样本对照整理、报告排版校对、模型表现汇总及汇报演示材料制作\\
\bottomrule
\end{tabularx}
\end{table}
